%%%%%%%%%%%%%%%%
%%% num 20 %%%
%%%%%%%%%%%%%%%%

\Chapter{20}

\Verse{1}
{
	\hR{וַיָּבֹאוּ בְנֵי יִשְׂרָאֵל כּל הָעֵדָה}
	\hP{מִדְבַּר צִן בַּחֹדֶשׁ הָרִאשׁוֹן וַיֵּשֶׁב הָעָם בְּקָדֵשׁ וַתָּמת שָׁם מִרְיָם וַתִּקָּבֵר שָׁם׃}
}
{
	\eR{And the children of Israel, all the congregation, came}
	\eP{to the wilderness of Zin in the first month, and the people
		stayed in Kadesh. And Miriam died there and was buried there.}
}
{
}

\Verse{2}
{
	\hP{וְלֹא הָיָה מַיִם לָעֵדָה וַיִּקָּהֲלוּ עַל מֹשֶׁה וְעַל אַהֲרֹן׃}
}
{
	\eP{
		And there was no water for the congregation, and they
		assembled against Moses and against Aaron.
	}
}
{
}

\Verse{3}
{
	\hP{וַיָּרֶב הָעָם עִם מֹשֶׁה וַיֹּאמְרוּ לֵאמֹר וְלוּ גָוַעְנוּ בִּגְוַע אַחֵינוּ לִפְנֵי יְהֹוָה׃}
}
{
	\eP{
		And the people quarreled with Moses, and they said,
		saying, “If only we had expired when our brothers expired
		in front of YHWH!
	}
}
{
}

\Verse{4}
{
	\hP{וְלָמָה הֲבֵאתֶם אֶת קְהַל יְהֹוָה אֶל הַמִּדְבָּר הַזֶּה לָמוּת שָׁם אֲנַחְנוּ וּבְעִירֵנוּ׃}
}
{
	\eP{
		And why have you brought YHWH’s community to this
		wilderness to die there, we and our cattle?
	}
}
{
}

\Verse{5}
{
	\hP{וְלָמָה הֶעֱלִיתֻנוּ מִמִּצְרַיִם לְהָבִיא אֹתָנוּ אֶל הַמָּקוֹם הָרָע הַזֶּה לֹא מְקוֹם זֶרַע וּתְאֵנָה וְגֶפֶן וְרִמּוֹן וּמַיִם אַיִן לִשְׁתּוֹת׃}
}
{
	\eP{
		And why did you bring us up from Egypt to bring us to this
		bad place? It’s not a place of seed and fig and vine and
		pomegranate, and there’s no water to drink!”
	}
}
{
}

\Verse{6}
{
	\hP{וַיָּבֹא מֹשֶׁה וְאַהֲרֹן מִפְּנֵי הַקָּהָל אֶל פֶּתַח אֹהֶל מוֹעֵד וַיִּפְּלוּ עַל פְּנֵיהֶם וַיֵּרָא כְבוֹד יְהֹוָה אֲלֵיהֶם׃}
}
{
	\eP{
		And Moses and Aaron came from in front of the community to
		the entrance of the Tent of Meeting, and they fell on
		their faces. And YHWH’s glory appeared to them.
	}
}
{
}

\Verse{7}
{
	\hP{וַיְדַבֵּר יְהֹוָה אֶל מֹשֶׁה לֵּאמֹר׃}
}
{
	\eP{And YHWH spoke to Moses, saying,}
}
{
}

\Verse{8}
{
	\hP{קַח אֶת הַמַּטֶּה וְהַקְהֵל אֶת הָעֵדָה אַתָּה וְאַהֲרֹן אָחִיךָ וְדִבַּרְתֶּם אֶל הַסֶּלַע לְעֵינֵיהֶם וְנָתַן מֵימָיו וְהוֹצֵאתָ לָהֶם מַיִם מִן הַסֶּלַע וְהִשְׁקִיתָ אֶת הָעֵדָה וְאֶת בְּעִירָם׃}
}
{
	\eP{
		“Take the staff and assemble the congregation, you and
		Aaron, your brother. And you shall speak to the rock
		before their eyes, and it will give its water. So you
		shall bring water out of the rock for them and give a
		drink to the congregation and their cattle.”
	}
}
{
%	\footnote{
%		Take the staff and … speak to the rock. Why does he need the staff if he
%		is only supposed to speak to the rock?! Some might say that God is
%		testing Moses. But the reason for taking the staff is already given in
%		the story of Aaron’s blossoming staff. Moses is told there to “put back
%		Aaron’s staff in front of the Testimony … for a sign to rebels” (Num
%		17:25). The text now says that Moses “took the staff from in front of
%		YHWH” (20:9). This expression would normally be expected to mean that he
%		took a staff that was located at the ark in the Tent of Meeting, “in
%		front of YHWH.” William Propp has pointed out that this connects back to
%		the staff of Aaron that miraculously blossomed and was placed before the
%		ark in the preceding episode. As quoted above, its purpose, explicitly,
%		was to be “a sign to rebels.” That is why Moses is supposed to carry it
%		in his hand in the people’s sight while dealing with their rebellion. As
%		he holds it, his first words to the people are in fact “Listen, rebels”
%		(19:10). And that is why Aaron is implicated by its misuse even though
%		Aaron does not say or do anything at the rock. It is his staff, with his
%		name on it, to be kept in a holy place and used in specific
%		circumstances. As we have seen in Leviticus, intent is not the issue in
%		the realm of ritual. Rather, actions relating to sacred objects and
%		boundaries bring necessary consequences of themselves. Whether or not
%		Aaron has thought or done anything improper, he is tied to what has
%		happened. Footnote 20:8. you shall speak. It is a plural: both Moses and
%		Aaron are supposed to speak to the rock. Footnote 20:8. rock. This is a
%		different word from the one in the other story of Moses striking a rock
%		to get water (Exod 17:1–7). There it is a crag (ûr) of Mount
%		Horeb (Sinai). Here it is a rock (sela‘) in the wilderness. In Moses’
%		last blessings to the people before he dies, he will speak of God as
%		Israel’s ûr (Deut 32:4,15,18,30,31,37); i.e., he uses the word
%		that recalls the water-from-the-rock episode that came out well, not the
%		present episode, which comes out disastrously for himself and Aaron.
%	}
}

\Verse{9}
{
	\hP{וַיִּקַּח מֹשֶׁה אֶת הַמַּטֶּה מִלִּפְנֵי יְהֹוָה כַּאֲשֶׁר צִוָּהוּ׃}
}
{
	\eP{
		And Moses took the staff from in front of YHWH as He
		commanded him.
	}
}
{
}

\Verse{10}
{
	\hP{וַיַּקְהִלוּ מֹשֶׁה וְאַהֲרֹן אֶת הַקָּהָל אֶל פְּנֵי הַסָּלַע וַיֹּאמֶר לָהֶם שִׁמְעוּ נָא הַמֹּרִים הֲמִן הַסֶּלַע הַזֶּה נוֹצִיא לָכֶם מָיִם׃}
}
{
	\eP{
		And Moses and Aaron assembled the community opposite the
		rock. And he said to them, “Listen, rebels, shall we bring
		water out of this rock for you?”
	}
}
{
}

\Verse{11}
{
	\hP{וַיָּרֶם מֹשֶׁה אֶת יָדוֹ וַיַּךְ אֶת הַסֶּלַע בְּמַטֵּהוּ פַּעֲמָיִם וַיֵּצְאוּ מַיִם רַבִּים וַתֵּשְׁתְּ הָעֵדָה וּבְעִירָם׃}
}
{
	\eP{
		And Moses lifted his hand and struck the rock with his
		staff twice. And much water came out! And the congregation
		and their cattle drank.
	}
}
{
%	\footnote{
%		water came out! Moses is supposed to speak to the rock, but he hits it.
%		So why does the miracle still work?! We might surmise that it means that
%		God makes it happen so as not to humiliate Moses. Or the explanation may
%		be that the miraculous power is in the staff, which is the one that
%		miraculously blossomed. Whatever the reason, though, the fact remains:
%		for the first time in the Bible, a human has changed a miracle. This is
%		an all-important step in a gradual shift in the balance of control of
%		miraculous phenomena in the Tanak. Adam, Noah, Abraham, and Isaac
%		perform no miracles. But Jacob manipulates the coloring of Laban’s
%		flock. Then Joseph interprets dreams. Then Moses and Aaron perform God’s
%		miracles. And now Moses changes a miracle. This shift will continue in
%		the biblical books that follow the Torah, and it is one of the central
%		developments of the Bible: Joshua will call for the sun to stand still
%		in the skies. By calling for a miracle on his own, without direction
%		from God, he goes even further than Moses. Later, Samson has powers
%		implanted in him at birth, so that he is free to use them as he wishes
%		all his life. Later still, Elijah and Elisha use miracles for a variety
%		of personal purposes. It appears that, starting with Moses, God is
%		entrusting humans with ever more responsibility and control of their
%		destiny. By the end of the Tanak, miracles cease. There are no visible
%		miracles in the late books of Ezra, Nehemiah, and Esther. Humans rather
%		must direct their destiny with their own natural human powers, not with
%		endowments of miraculous powers from God. Through the course of the
%		Tanak, humans are forced to grow up and become more self-reliant.
%	}
}

\Verse{12}
{
	\hP{וַיֹּאמֶר יְהֹוָה אֶל מֹשֶׁה וְאֶל אַהֲרֹן יַעַן לֹא הֶאֱמַנְתֶּם בִּי לְהַקְדִּישֵׁנִי לְעֵינֵי בְּנֵי יִשְׂרָאֵל לָכֵן לֹא תָבִיאוּ אֶת הַקָּהָל הַזֶּה אֶל הָאָרֶץ אֲשֶׁר נָתַתִּי לָהֶם׃}
}
{
	\eP{
		And YHWH said to Moses and to Aaron, “Because you did not
		trust in me, to make me holy before the eyes of the
		children of Israel, therefore you shall not bring this
		community to the land that I have given them!”
	}
}
{
%	\footnote{
%		you did not trust in me, to make me holy. What is Moses’ offense? What
%		has he done that is so horrible as to deserve what is probably the worst
%		possible thing that could have been done to him: that he will not live
%		to enter the land? Many answers have been proposed to the question of
%		what the sin of Moses is. He calls the people rebels. He says, “Shall we
%		bring water out of this rock” instead of “Shall God … ” He strikes the
%		rock instead of speaking to it. The answer may lie in the wording of
%		YHWH’s reprimand: “You did not trust in me, to make me holy.” By hitting
%		the rock Moses makes it appear more as if it is his own agency that is
%		producing the water than if he had just spoken to the rock and watched
%		the water flow out. His saying “we” has the same effect verbally that
%		his physical action has visually. True, YHWH had said to him in his
%		instructions, “So you shall bring water out,” but Moses’ saying it to
%		the people before hitting the rock still makes those words appear to
%		mean something quite different from their meaning in the original
%		instructions. By word and act Moses is thus appropriating to himself an
%		act of God. In doing this he is undoing the message that God and Moses
%		himself have been conveying to the people up to this point. The people
%		have continuously directed their attention to Moses instead of to God.
%		In this story, as in the others so far, they say to Moses and Aaron,
%		“Why have you brought YHWH’s community to this wilderness … ? Why did
%		you bring us up from Egypt?” Until this episode Moses has repeatedly
%		told the people, “It is not from my own heart” and “You are congregating
%		against YHWH,” but now his words and actions confirm the people’s own
%		perception. He has come a considerable way from his timid attempts to
%		avoid the assignment at the burning bush. At Meribah Moses oversteps a
%		divine limit, and he pays a maximal price. The severity of the
%		punishment demonstrates the seriousness of the offense. Footnote 20:12.
%		you shall not bring this community to the land. In a way, the sin of
%		Moses is like the sin of Adam and Eve: stepping over the line, taking a
%		power from the divine realm. His punishment is not death. He lives to be
%		120 years old, the maximum (as decreed in Gen 6:3); his punishment is
%		exile—he is not permitted to enter the land—which is also like the
%		punishment of Adam and Eve, who are prevented from reentering Eden.
%	}
}

\Verse{13}
{
	\hP{הֵמָּה מֵי מְרִיבָה אֲשֶׁר רָבוּ בְנֵי יִשְׂרָאֵל אֶת יְהֹוָה וַיִּקָּדֵשׁ בָּם׃}
}
{
	\eP{
		They are the waters of Meribah, over which the children of
		Israel quarreled with YHWH, and He was made holy among
		them.
	}
}
{
}

\Verse{14}
{
	\hJ{וַיִּשְׁלַח מֹשֶׁה מַלְאָכִים מִקָּדֵשׁ אֶל מֶלֶךְ אֱדוֹם כֹּה אָמַר אָחִיךָ יִשְׂרָאֵל אַתָּה יָדַעְתָּ אֵת כּל הַתְּלָאָה אֲשֶׁר מְצָאָתְנוּ׃}
}
{
	\eJ{
		And Moses sent messengers from Kadesh to the king of Edom.
		“Your brother, Israel, says this: You know all the
		hardship that has found us:
	}
}
{
}

\Verse{15}
{
	\hJ{וַיֵּרְדוּ אֲבֹתֵינוּ מִצְרַיְמָה וַנֵּשֶׁב בְּמִצְרַיִם יָמִים רַבִּים וַיָּרֵעוּ לָנוּ מִצְרַיִם וְלַאֲבֹתֵינוּ׃}
}
{
	\eJ{
		And our fathers went down to Egypt, and we lived in Egypt
		many days. And Egypt was bad to us and to our fathers,
	}
}
{
}

\Verse{16}
{
	\hJ{וַנִּצְעַק אֶל יְהֹוָה וַיִּשְׁמַע קֹלֵנוּ וַיִּשְׁלַח מַלְאָךְ וַיֹּצִאֵנוּ מִמִּצְרָיִם וְהִנֵּה אֲנַחְנוּ בְקָדֵשׁ עִיר קְצֵה גְבוּלֶךָ׃}
}
{
	\eJ{
		and we cried to YHWH, and He heard our voice and sent an
		angel and brought us from Egypt. And here we are in
		Kadesh, a city at the edge of your border.
	}
}
{
}

\Verse{17}
{
	\hJ{נַעְבְּרָה נָּא בְאַרְצֶךָ לֹא נַעֲבֹר בְּשָׂדֶה וּבְכֶרֶם וְלֹא נִשְׁתֶּה מֵי בְאֵר דֶּרֶךְ הַמֶּלֶךְ נֵלֵךְ לֹא נִטֶּה יָמִין וּשְׂמֹאול עַד אֲשֶׁר נַעֲבֹר גְּבֻלֶךָ׃}
}
{
	\eJ{
		Let us pass through your land. We won’t pass through a
		field or through a vineyard, and we won’t drink the water
		of a well. We’ll go by the king’s road, we won’t turn
		right or left, until we pass your border.”
	}
}
{
}

\Verse{18}
{
	\hJ{וַיֹּאמֶר אֵלָיו אֱדוֹם לֹא תַעֲבֹר בִּי פֶּן בַּחֶרֶב אֵצֵא לִקְרָאתֶךָ׃}
}
{
	\eJ{
		And Edom said to him, “You won’t pass through me, or else
		I’ll go out to you with the sword.”
	}
}
{
}

\Verse{19}
{
	\hJ{וַיֹּאמְרוּ אֵלָיו בְּנֵי יִשְׂרָאֵל בַּמְסִלָּה נַעֲלֶה וְאִם מֵימֶיךָ נִשְׁתֶּה אֲנִי וּמִקְנַי וְנָתַתִּי מִכְרָם רַק אֵין דָּבָר בְּרַגְלַי אֶעֱבֹרָה׃}
}
{
	\eJ{
		And the children of Israel said to him, “We’ll go up by
		the highway, and if we drink your water, I and my cattle,
		then I’ll pay its price. Only, it’s nothing: let me pass
		through by foot.”
	}
}
{
}

\Verse{20}
{
	\hJ{וַיֹּאמֶר לֹא תַעֲבֹר וַיֵּצֵא אֱדוֹם לִקְרָאתוֹ בְּעַם כָּבֵד וּבְיָד חֲזָקָה׃}
}
{
	\eJ{
		And he said, “You won’t pass.” And Edom went out to him
		with a heavy mass of people and a strong hand,
	}
}
{
}

\Verse{21}
{
	\hJ{וַיְמָאֵן אֱדוֹם נְתֹן אֶת יִשְׂרָאֵל עֲבֹר בִּגְבֻלוֹ וַיֵּט יִשְׂרָאֵל מֵעָלָיו׃}
}
{
	\eJ{
		and Edom refused to allow Israel to pass through his
		border. And Israel turned away from him.
	}
}
{
}

\Verse{22}
{
	\hR{וַיִּסְעוּ מִקָּדֵשׁ וַיָּבֹאוּ בְנֵי יִשְׂרָאֵל כּל הָעֵדָה הֹר הָהָר׃}
}
{
	\eR{
		And they traveled from Kadesh, and the children of Israel,
		all the congregation, came to Mount Hor.
	}
}
{
}

\Verse{23}
{
	\hP{וַיֹּאמֶר יְהֹוָה אֶל מֹשֶׁה וְאֶל אַהֲרֹן בְּהֹר הָהָר עַל גְּבוּל אֶרֶץ אֱדוֹם לֵאמֹר׃}
}
{
	\eP{
		And YHWH said to Moses and to Aaron at Mount Hor, on the
		border of the land of Edom, saying,
	}
}
{
}

\Verse{24}
{
	\hP{יֵאָסֵף אַהֲרֹן אֶל עַמָּיו כִּי לֹא יָבֹא אֶל הָאָרֶץ אֲשֶׁר נָתַתִּי לִבְנֵי יִשְׂרָאֵל עַל אֲשֶׁר מְרִיתֶם אֶת פִּי לְמֵי מְרִיבָה׃}
}
{
	\eP{
		“Let Aaron be gathered to his people, because he shall not
		come to the land that I have given to the children of
		Israel, because you rebelled against my word at the waters
		of Meribah.
	}
}
{
%	\footnote{
%		you rebelled. You rebelled—meaning Moses and Aaron! God tells them that
%		they themselves did exactly what they were supposed to stop the people
%		from doing. It is especially painful for Moses, who said “Listen,
%		rebels,” to hear his God apply that word now to him. Leaders of a
%		congregation cannot violate the very instruction that they uphold and
%		teach to others.
%	}
}

\Verse{25}
{
	\hP{קַח אֶת אַהֲרֹן וְאֶת אֶלְעָזָר בְּנוֹ וְהַעַל אֹתָם הֹר הָהָר׃}
}
{
	\eP{
		Take Aaron and Eleazar, his son, and take them up Mount
		Hor,
	}
}
{
}

\Verse{26}
{
	\hP{וְהַפְשֵׁט אֶת אַהֲרֹן אֶת בְּגָדָיו וְהִלְבַּשְׁתָּם אֶת אֶלְעָזָר בְּנוֹ וְאַהֲרֹן יֵאָסֵף וּמֵת שָׁם׃}
}
{
	\eP{
		and take off Aaron’s clothes, and you shall put them on
		Eleazar, his son. And Aaron will be gathered and die
		there.”
	}
}
{
}

\Verse{27}
{
	\hP{וַיַּעַשׂ מֹשֶׁה כַּאֲשֶׁר צִוָּה יְהֹוָה וַיַּעֲלוּ אֶל הֹר הָהָר לְעֵינֵי כּל הָעֵדָה׃}
}
{
	\eP{
		And Moses did as YHWH commanded him. And they went up to
		Mount Hor before the eyes of all the congregation.
	}
}
{
}

\Verse{28}
{
	\hP{וַיַּפְשֵׁט מֹשֶׁה אֶת אַהֲרֹן אֶת בְּגָדָיו וַיַּלְבֵּשׁ אֹתָם אֶת אֶלְעָזָר בְּנוֹ וַיָּמת אַהֲרֹן שָׁם בְּרֹאשׁ הָהָר וַיֵּרֶד מֹשֶׁה וְאֶלְעָזָר מִן הָהָר׃}
}
{
	\eP{
		And Moses took off Aaron’s clothes and put them on
		Eleazar, his son. And Aaron died there on the top of the
		mountain. And Moses and Eleazar came down from the
		mountain.
	}
}
{
}

\Verse{29}
{
	\hP{וַיִּרְאוּ כּל הָעֵדָה כִּי גָוַע אַהֲרֹן וַיִּבְכּוּ אֶת אַהֲרֹן שְׁלֹשִׁים יוֹם כֹּל בֵּית יִשְׂרָאֵל׃}
}
{
	\eP{
		And all the congregation saw that Aaron had expired, and
		all the house of Israel mourned Aaron thirty days.
	}
}
{
}
